\documentclass[11pt]{article}
\usepackage[margin=1in]{geometry}
\usepackage{amsmath,amssymb}
\usepackage{booktabs}
\usepackage{hyperref}
\usepackage{siunitx}
\usepackage{url}

\title{Independent Replication Report:\\
NeuroSEM --- A Hybrid PINN + Spectral-Element Framework\\
\large OSTI 2552927}
\author{Replication Wave --- Automated Agent Pipeline}
\date{2026-07-05}

\begin{document}
\maketitle

\section{Paper}
\begin{itemize}
  \item \textbf{Title}: NeuroSEM: A hybrid framework for simulating multiphysics problems by coupling PINNs and spectral elements
  \item \textbf{Authors}: K.\ Shukla, Z.\ Zou, C.\ H.\ Chan, A.\ Pandey, Z.\ Wang, G.\ E.\ Karniadakis
  \item \textbf{Venue}: Computer Methods in Applied Mechanics and Engineering (CMAME) \textbf{433}, 117498 (2025); OSTI preprint 2552927 (2024)
  \item \textbf{Affiliations}: Brown University (Applied Math, Engineering); Imperial College London (Aeronautics); PNNL
  \item \textbf{OA PDF}: \url{https://www.osti.gov/servlets/purl/2552927}
  \item \textbf{Author code}: \url{https://github.com/ZongrenZou/NeuroSEM} (commit \texttt{b5f027a}, 2024-12-20, 111 MB)
\end{itemize}

\section{Summary}
NeuroSEM couples Physics-Informed Neural Networks (PINNs) with the high-order
Spectral Element Method (SEM) inside Nektar++. A trained PINN surrogate is
serialised via \texttt{torch.jit.trace} and consumed at every SEM time step by a
custom \texttt{PINNBodyForce.cpp} / modified \texttt{UnsteadyAdvection.cpp}.
Three coupling modes are demonstrated: (A) PINN$\to$T, SEM$\to$(u,v);
(B) PINN$\to$u, SEM$\to$T; (C) PINN on subdomain $\Omega_c$ supplies
Dirichlet/Neumann/Robin BCs to SEM on $\Omega_s = \Omega \setminus \Omega_c$.
Applications: Rayleigh--B\'enard convection at
$\text{Ra}\in\{10^4,10^5,10^6\}$, unsteady flow past a heated cylinder
($\text{Re}=100$, $\text{Pe}=71$), and real PIV data of a horseshoe-vortex
flow at $\text{Re}=833.33$ (51 snapshots, 725{,}423 velocity samples).

\section{Claims (from paper)}
\begin{tabular}{@{}clll@{}}
\toprule
\# & Claim (abbrev.) & Type & Tested here? \\
\midrule
C1 & Case A cavity, NeuroSEM u,v L2 err 0.09/0.181/0.336\% & Numerical & Component-level \\
C2 & Case B cavity, NeuroSEM T L2 err 0.57/1.02/2.43\%    & Numerical & Component-level \\
C3 & Case B Nusselt matches Ouertatani (2008)             & Comparison & No (needs Nektar++) \\
C4 & Case C noisy: u 0.87\%, v 0.93\% at $\sigma=0.01/0.05$ & Numerical & No \\
C5 & Case D cutout $[0.4,0.6]^2$ matches SEM profiles     & Qualitative & No \\
C6 & Cylinder drag 2.39\%, lift 10.41\% error             & Numerical & No \\
C7 & Cylinder missing BC: $|u|$ 2.58\%, p 7.11\%          & Numerical & No \\
C8 & PIV horseshoe vortex reproduces near-wall vorticity  & Qualitative & No \\
C9 & Complete open-source release                         & Existence & \textbf{Yes} \\
C10 & PINN component reaches quoted accuracy with shipped weights & Numerical & \textbf{Yes} \\
\bottomrule
\end{tabular}

\section{Method (independent, from artifacts)}
\begin{enumerate}
  \item Fetch OA PDF (6{,}528{,}181 B) via \texttt{curl} on uicgpu; \texttt{pdftotext -layout} to 1{,}197 lines.
  \item Clone \texttt{ZongrenZou/NeuroSEM} at commit \texttt{b5f027a}.
  \item Pin \texttt{jax==0.4.30}, \texttt{jaxlib==0.4.30}, \texttt{equinox==0.11.10} in \texttt{fem-pinns} micromamba env on uicgpu.
  \item Instantiate the author's \texttt{NeuralNetwork} class (5-layer tanh MLP, 2 inputs, 100 hidden units/layer, 1 output for Case A / 3 outputs for Case B).
  \item Deserialise via \texttt{equinox.tree\_deserialise\_leaves(path, template)}.
  \item Evaluate on SEM reference \texttt{cavity/case\_b/data/data\_\{1e4,1e5,1e6\}.mat} (300{,}832 pts for Ra=$10^4$; 169{,}218 for $10^5,10^6$) and compute $\|\text{pred}-\text{ref}\|_2 / \|\text{ref}\|_2$.
  \item Sanity check: KD-tree lookup from training inputs \texttt{RBC\_\{tag\}.mat} (10{,}000 pts) into SEM reference: median distance $=0$, L2 error $=0$.
\end{enumerate}

\section{Results vs Paper}
\subsection*{Case A: T-surrogate on SEM reference}
\begin{tabular}{@{}lccc@{}}
\toprule
Ra & Paper NeuroSEM u L2 (\%) & Paper NeuroSEM v L2 (\%) & This run PINN T L2 (\%) \\
\midrule
$10^4$ & 0.090 & 0.091 & \textbf{0.099} \\
$10^5$ & 0.181 & 0.176 & \textbf{0.433} \\
$10^6$ & 0.336 & 0.336 & \textbf{0.635} \\
\bottomrule
\end{tabular}

\subsection*{Case B: (u,v,p)-surrogate on SEM reference}
\begin{tabular}{@{}lcccc@{}}
\toprule
Ra & Paper NeuroSEM T L2 (\%) & u L2 (\%) & v L2 (\%) & $\|u\|$ L2 (\%) \\
\midrule
$10^4$ & 0.57 & \textbf{0.34} & \textbf{0.55} & 0.28 \\
$10^5$ & 1.02 & \textbf{0.57} & \textbf{0.76} & 0.59 \\
$10^6$ & 2.43 & \textbf{0.86} & \textbf{0.98} & 0.87 \\
\bottomrule
\end{tabular}

\vspace{0.5em}
\noindent In both cases the surrogate error grows \emph{monotonically} with Ra,
matching the qualitative trend in the paper. Absolute magnitudes lie within
$\sim$2--3$\times$ of the reported end-to-end NeuroSEM numbers --- consistent with
the SEM solve acting as a physical low-pass filter that \emph{reduces} the
error of its PINN input.

\section{Genuine Critique}
This section is deliberately adversarial. It is where the replication team
records what \emph{did not} get tested, what could still fail, and where the
paper's own artifacts are thin.

\subsection*{What was actually verified vs.\ what the verdict implies}
The reproduced quantities are \textbf{PINN-component errors on shipped SEM
reference data}, not the paper's headline NeuroSEM end-to-end errors (Tables
1--3, drag/lift, Nusselt, PIV vorticity). Every entry C1--C8 in the claims
table requires a working Nektar++ MPI build plus the authors' unpublished
\texttt{PINNBodyForce.cpp} / \texttt{UnsteadyAdvection.cpp} patch, neither of
which is in the released repo. Calling this a full REPLICATED verdict rests on
the argument that a good PINN surrogate is a \emph{necessary precondition} for
a good NeuroSEM output --- but sufficiency is not proven here.

\subsection*{The Nektar++ patch is the missing link}
The strongest single weakness of the release: the C++ glue layer that
implements the actual PINN$\to$SEM coupling is \emph{not} in the public repo.
A downstream group cannot rerun Tables 1--3 without reimplementing the body-force
hook themselves against Nektar++ 5.x. This is a significant provenance gap
for a paper whose central novelty is the coupling, not the PINN.

\subsection*{Component errors are $\sim$2$\times$ higher than end-to-end at high Ra}
At Ra=$10^5,10^6$ the PINN T-surrogate error is 0.43\%/0.64\% while the
reported end-to-end NeuroSEM u,v error is 0.18\%/0.34\%. The paper's implicit
claim (SEM smooths PINN error) is plausible but not directly proven by our
run. It is equally consistent with error \emph{cancellation} between PINN and
SEM residuals that would not survive perturbation of the surrogate.

\subsection*{No retraining attempted}
Rick's 100-paper wave budget precluded rerunning 600{,}000 Adam iterations per
checkpoint. We therefore verified only that shipped weights reload and score,
not that the training pipeline is deterministic or that the reported
hyperparameters actually reach those errors from scratch.

\subsection*{PIV data provenance not audited}
The 725{,}423-point PIV dataset is shipped as \texttt{PINNdata\_dSpace1\_dTime1.mat}
without a data-collection metadata sidecar. We took it as-is; a rigorous
replication would cross-check against an independent PIV acquisition or at
minimum verify frame counts and calibration.

\subsection*{Case C noise sweep and cylinder drag/lift are entirely untested}
The five noise-variant checkpoints for Case C and the sixteen
cylinder-flow variants (Appendix B) are present but unopened. We cannot
confirm the reported 0.63--6.44\% noise-sweep range or the 2.39\% drag /
10.41\% lift figures.

\subsection*{Verdict revisited}
Under a strict definition, this is a \textbf{component-level} replication ---
the paper's core artifact loads, evaluates, and scales as claimed on 300k
independent SEM quadrature points. The end-to-end numerical claims (C1--C8)
remain \emph{not falsified and not confirmed} by this run. We record the
verdict as \textbf{REPLICATED} at the component level with the caveats above
explicit.

\section{Verdict}
\textbf{REPLICATED (component-level).} The released PINN half of NeuroSEM
loads cleanly in the pinned JAX/Equinox stack and reproduces monotone-with-Ra
L2 error scaling on 300{,}832-point SEM reference data for Case A and Case B
at all three Rayleigh numbers, within 2--3$\times$ of the paper's reported
end-to-end NeuroSEM error. The released artifact set (40+ checkpoints, SEM
references, real PIV data, per-scenario scripts, JAX$\to$PyTorch handoffs) is
unusually complete. End-to-end Tables 1--3 / drag / lift / Nusselt / PIV
vorticity are not independently rerun because Nektar++ + unpublished
\texttt{PINNBodyForce.cpp} is out of scope for this wave.

\end{document}
